\documentclass[10pt,a4paper]{article}
\title{Techniques de Simulations}
\author{
  PRISCILLA, GOGUY\\
  \texttt{priscilla.goguy@etu.univ-lyon1.fr}
  \and
  MAHLÎ, REINETTE\\
  \texttt{mahli.reinette@etu.univ-lyon1.fr}
}
\date{13/10/2025}
\usepackage[utf8]{inputenc}
\usepackage[T1]{fontenc}
\usepackage[french]{babel}
\usepackage{graphicx}
\usepackage{hyperref}

\usepackage{listings}
\usepackage{xcolor}


\usepackage[many]{tcolorbox}
\usepackage{lipsum}
\usepackage{tikz}
\usetikzlibrary{automata, positioning, arrows}
\usepackage{pgfplots}
\usepackage{xcolor}


% le préambule











% initialisation des couleurs
\definecolor{codegreen}{rgb}{0,0.6,0}
\definecolor{codegray}{rgb}{0.5,0.5,0.5}
\definecolor{codepurple}{rgb}{0.58,0,0.82}
\definecolor{backcolour}{rgb}{0.95,0.95,0.92}






%block de code 
\lstdefinestyle{mystyle}{
    backgroundcolor=\color{backcolour},   
    commentstyle=\color{codegreen},
    keywordstyle=\color{magenta},
    numberstyle=\tiny\color{codegray},
    stringstyle=\color{codepurple},
    basicstyle=\ttfamily\footnotesize,
    breakatwhitespace=false,         
    breaklines=true,                 
    captionpos=b,                    
    keepspaces=true,                 
    numbers=left,                    
    numbersep=5pt,                  
    showspaces=false,                
    showstringspaces=false,
    showtabs=false,                  
    tabsize=2
}

\lstset{style=mystyle}






%initialisation graphics
\tikzset{
->,
node distance = 2cm}





















% le corps du document






% le préambule
\begin{document}
%titre
\pagecolor{black!20}
\maketitle

\begin{center}
\includegraphics[width=4cm,height=2cm]{img1}

\href{https://github.com/LaboiteNoire/techniques-de-simulations-}{\includegraphics[width=1cm,height=1cm]{img2}}

\end{center}


\newpage

\phantom{aaaaaa}

\tableofcontents

\part{Préambule}
\section{Avant-propos}
\section{Notations}

\part{Modélisation}


\part{Conclusion}
\section{Commentaires}
\section{Annexes et contacts}






\end{document}